\documentclass[12pt]{article}
\usepackage[T1]{fontenc}
\usepackage[utf8]{inputenc}
\usepackage{lmodern}
\usepackage{textcomp}
\usepackage{enumitem}
\usepackage{geometry}
\geometry{margin=1in}
\begin{document}
\begin{center}
Answer Key - Computer Science - K-12 Level Assessment 22 \
\small (Answers are ordered exactly as in the question paper.)
\end{center}

\section*{Multiple Choice}
\begin{enumerate}[label=\arabic*.]
\item \textbf{Answer:} D
\\ \textit{Options:} A) A function that averages numeric values in a column or row; B) A function that counts the values in a column or row; C) A function that rounds a numeric value; D) A function that sorts values in a column or row
\\ \textit{Note:} Sorting values in a column or row allows users to easily identify outliers or anomalies in the data set, which are often indicative of data entry errors.
\item \textbf{Answer:} D
\\ \textit{Options:} A) 1; B) 2; C) 3; D) 4
\\ \textit{Note:} The correct answer is 4 because the `max()` function in Python 3 returns the largest element from a list, and in the list l = [1, 2, 3, 4], the largest number is 4.
\item \textbf{Answer:} A
\\ \textit{Options:} A) 4; B) 1; C) 256; D) 16
\\ \textit{Note:} The correct answer is 4 because in Python, the expression is evaluated as 4 multiplied by 1 raised to the power of 4, and since any number raised to the power of 0 is 1, the calculation simplifies to 4 * 1, which equals 4.
\item \textbf{Answer:} C
\\ \textit{Options:} A) [ 'abcd', 786 , 2.23, 'john', 70.2 ]; B) abcd; C) [786, 2.23]; D) None of the above.
\\ \textit{Note:} The correct answer is [786, 2.23] because in Python, slicing a list with the notation list[start:end] retrieves the elements from index 'start' up to, but not including, index 'end', so list[1:3] returns the elements at indices 1 and 2.
\item \textbf{Answer:} A
\\ \textit{Options:} A) a[i] == max; B) a[i] != max; C) a[i] \textless{} max || a[i] \textgreater{} max; D) FALSE
\\ \textit{Note:} The expression !(max != a[i]) is equivalent to saying that max is equal to a[i], thus simplifying the entire boolean expression to a[i] == max.
\end{enumerate}

\section*{True / False}
\begin{enumerate}[label=\arabic*.]
\setcounter{enumi}{5}
\item \textbf{Answer:} False
\\ \textit{Options:} A) True; B) False
\\ \textit{Note:} In Python 3, the function random() generates a random float between 0.0 and 1.0 but does not set or alter the integer starting value for generating random numbers, which is done using the seed() function instead.
\item \textbf{Answer:} True
\\ \textit{Options:} A) True; B) False
\\ \textit{Note:} The statement is true because the lambda function defined as r = lambda x: 2 * x will return 2 times the input value, so r(3) evaluates to 2 * 3, which equals 6.
\item \textbf{Answer:} True
\\ \textit{Options:} A) True; B) False
\\ \textit{Note:} Authentication is compromised because the use of a stolen login and password allows unauthorized individuals to gain access to the system, bypassing the intended security measures.
\item \textbf{Answer:} True
\\ \textit{Options:} A) True; B) False
\\ \textit{Note:} In Python 3, the `set` function creates a collection of unique elements, automatically removing any duplicate values from the input list, resulting in the output \{1, 2, 3, 4\}.
\item \textbf{Answer:} False
\\ \textit{Options:} A) True; B) False
\\ \textit{Note:} The statement is false because the condition a \textless{} b alone does not account for other factors or contexts that might influence the truth value of the expression, such as data types or additional logical conditions that may be involved.
\end{enumerate}

\section*{Long Form Response}
\begin{enumerate}[label=\arabic*.]
\setcounter{enumi}{10}
\item \textbf{Answer:} def sqrt\_root(num): | return sqrt\_root
\\ \textit{Note:} The answer correctly defines a function named `sqrt\_root` that takes a numerical input, indicating an intent to find the square root of a perfect number, though it lacks the implementation detail for the actual calculation.
\item \textbf{Answer:} def get\_median(arr 1, arr 2, n): | return (m 1 + m 2)/2
\\ \textit{Note:} correct because it accurately describes the process of calculating the median by averaging the two middle values (m 1 and m 2) of the combined sorted arrays, assuming that the function properly identifies those values based on the input arrays.
\end{enumerate}

\vfill
\noindent\textit{End of Answer Key}
\end{document}